\section{Discusión}
\label{sec:discusion}

Los resultados muestran que la especificación formal captura
correctamente las propiedades centrales del algoritmo de búsqueda. Las
postcondiciones de \texttt{buscarRutasDirectas} y
\texttt{buscarRutasCombinadas} expresan con precisión las garantías del
sistema: proximidad de las rutas al origen y al destino, y viabilidad de
los transbordos. La ejecución de los casos de prueba y la evaluación
directa de la postcondición confirman que el modelo satisface dichas
propiedades para los datos considerados.

Sin embargo, la validación por interpretación no equivale a una
demostración deductiva completa: los casos de prueba cubren escenarios
representativos, pero no prueban la postcondición para todas las
entradas posibles. Las obligaciones de prueba generadas por la Toolbox
constituyen el punto de partida para una verificación más rigurosa.

Desde el punto de vista del modelado, la herramienta impuso algunas
restricciones prácticas. En particular, los tipos registro se
representaron como tipos producto (tuplas) y la distancia se modeló de
forma simplificada; ambos aspectos se detallan en la
Sección~\ref{sec:limitaciones}.
